\documentclass[12pt,a4paper]{article}

% ----- kódolás, nyelv -----
\usepackage[utf8]{inputenc}
\usepackage[T1]{fontenc}
\usepackage[hungarian]{babel}

% ----- betűtípus (Times New Roman közelítő) -----
\usepackage{mathptmx}

% ----- oldalgeometria -----
\usepackage[top=2.5cm, bottom=2.5cm, left=3cm, right=2.5cm]{geometry}

% ----- sortávolság -----
\usepackage{setspace}
\setstretch{1.2}

% ----- bekezdések -----
\setlength{\parindent}{0pt}
\setlength{\parskip}{6pt}

% ----- matematika -----
\usepackage{amsmath}

% ----- színek -----
\usepackage[dvipsnames,table]{xcolor}

% ----- képek -----
\usepackage{graphicx}
\usepackage{float}

% ----- táblázatok -----
\usepackage{booktabs}
\usepackage{tabularx}

% ----- kódlisták -----
\usepackage{listings}

\lstdefinestyle{java}{
  language         = Java,
  basicstyle       = \ttfamily\small,
  keywordstyle     = \color{NavyBlue}\bfseries,
  commentstyle     = \color{OliveGreen}\itshape,
  stringstyle      = \color{BrickRed},
  numbers          = left,
  numberstyle      = \tiny\color{gray},
  numbersep        = 8pt,
  frame            = single,
  rulecolor        = \color{gray!40},
  breaklines       = true,
  tabsize          = 4,
  showstringspaces = false,
  literate =
    {á}{{\'{a}}}1 {é}{{\'{e}}}1 {í}{{\'\i}}1  {ó}{{\'{o}}}1
    {ö}{{\"o}}1   {o"}{{\H{o}}}1 {ő}{{\H{o}}}1 {ú}{{\'{u}}}1
    {ü}{{\"u}}1   {u"}{{\H{u}}}1 {ű}{{\H{u}}}1
    {Á}{{\'{A}}}1 {É}{{\'{E}}}1 {Í}{{\'{I}}}1  {Ó}{{\'{O}}}1
    {Ö}{{\"O}}1   {Ő}{{\H{O}}}1 {Ú}{{\'{U}}}1
    {Ü}{{\"U}}1   {Ű}{{\H{U}}}1,
}

\lstdefinestyle{shell}{
  basicstyle       = \ttfamily\small,
  frame            = single,
  rulecolor        = \color{gray!40},
  backgroundcolor  = \color{gray!7},
  breaklines       = true,
  numbers          = none,
}

\lstdefinestyle{plain}{
  basicstyle       = \ttfamily\small,
  frame            = single,
  rulecolor        = \color{gray!40},
  breaklines       = true,
  numbers          = none,
}

% ----- hivatkozások -----
\usepackage[hidelinks]{hyperref}

% ----- fejléc / lábléc -----
\usepackage{fancyhdr}
\pagestyle{fancy}
\fancyhf{}
\fancyhead[L]{\small OOP féléves feladat}
\fancyhead[R]{\small JavaFX 3D Teáskanna Nézegető}
\fancyfoot[C]{\thepage}
\renewcommand{\headrulewidth}{0.4pt}

% ----- inline kód -----
\newcommand{\code}[1]{\texttt{#1}}

% ----- képhely placeholder -----
% Használat: \kephely[magasság]{felirat}
\newcommand{\kephely}[2][6cm]{%
  \begin{figure}[H]
    \centering
    \setlength{\fboxsep}{14pt}%
    \fbox{%
      \begin{minipage}[c][#1][c]{0.88\textwidth}
        \centering
        {\large\color{gray}\itshape [Képernyőkép helye]}\\[0.6em]
        {\normalsize\color{gray!80}\bfseries #2}
      \end{minipage}%
    }
    \caption{#2}
  \end{figure}%
}

% ============================================================
\begin{document}
% ============================================================

% ---- CÍMLAP ----
\begin{titlepage}
  \centering
  \vspace*{1.8cm}

  {\large Miskolci Egyetem}\\[0.4cm]
  {\large Informatikai Intézet}\\[2.2cm]

  {\LARGE\bfseries Objektumorientált programozás}\\[0.6cm]
  {\Large Féléves feladat dokumentáció}\\[1.6cm]

  \rule{0.55\textwidth}{0.5pt}\\[1.4cm]

  {\Large\bfseries JavaFX 3D Teáskanna Nézegető}\\[2.4cm]

  \begin{tabular}{ll}
    \textbf{Hallgató neve:} & {[Hallgató neve]}               \\[0.35cm]
    \textbf{Neptun kód:}    & {[Neptun kód]}                  \\[0.35cm]
    \textbf{Csoport:}       & {[Csoport]}                     \\[0.35cm]
    \textbf{Oktató:}        & {[Oktató neve]}                 \\[0.35cm]
    \textbf{Tanév:}         & 2024/2025. tanév, tavaszi félév \\[0.35cm]
    \textbf{Leadás dátuma:} & {[Dátum]}                       \\
  \end{tabular}

  \vfill
  {\small Miskolc, 2025}
\end{titlepage}

% ---- TARTALOMJEGYZÉK ----
\tableofcontents
\newpage

% ============================================================
\section{Bevezetés}
% ============================================================

\subsection{A feladat ismertetése}

A feladat célja egy ismert, hétköznapi tárgy háromdimenziós modelljének elkészítése és
interaktív megjelenítése JavaFX segítségével. A választott tárgy egy teáskanna, amelynek
geometriai felépítése egyszerű forgástestek és összetett hálógenerálási technikák
kombinációjával közelíthető meg.

Az alkalmazás lehetővé teszi, hogy a felhasználó a teáskanna összes fontosabb geometriai
paraméterét valós időben állítsa be. A módosítások azonnal megjelennek a 3D nézetben:
a modell újragenerálódik, és a csúszkák értékeinek megfelelő alakot vesz fel. A programban
megtalálható a mentés és betöltés funkció is, így a beállítások JSON fájlba exportálhatók,
majd visszatölthetők.

\kephely[7.5cm]{A program főablaka -- teljes nézet}

\subsection{A program funkciói}

Az alkalmazás az alábbi fő funkciókat valósítja meg:

\begin{itemize}
  \item \textbf{Generálás:} Az \textit{Új} menüpont véletlenszerű paraméterekkel hoz létre
        egy új teáskannaformát.
  \item \textbf{Törlés / Visszaállítás:} A \textit{Törlés} menüpont az összes paramétert
        visszaállítja az alapértékekre.
  \item \textbf{Mentés (Exportálás):} Az aktuális paraméterek JSON formátumban
        fájlba menthetők.
  \item \textbf{Betöltés (Importálás):} Korábban elmentett JSON fájlból visszatölthetők
        a paraméterek.
  \item \textbf{Rajzolás / Megjelenítés:} A 3D modell valós időben frissül minden
        paraméterváltozásnál.
  \item \textbf{Bemutatás:} Az egérrel a modell szabadon forgatható, az egérgörgővel
        nagyítható és kicsinyíthető.
\end{itemize}

Az objektum öt alkatrészből áll: a váza test, a kiöntő, a fogantyú, a tetőfélgömb és
a tető fogógombja. Mindegyik alkatrész saját színnel és paraméterekkel rendelkezik.

\kephely[6cm]{A teáskanna öt alkatrésze -- jelölt nézet}

\subsection{Elméleti háttér}

A program középpontjában a JavaFX \code{TriangleMesh} típusa áll. A háromszöghálós
reprezentáció a 3D grafika egyik legszélesebb körben használt módszere: az objektum
felszínét egymáshoz illeszkedő háromszögek összessége közelíti. Minden háromszöget három
csúcspont (vertex) határoz meg, az egymás melletti háromszögek közös csúcsokat és éleket
osztanak meg.

A váza körvonalát egy profil-görbe írja le, amelyet forgástestként körbeforgatunk a
függőleges tengely körül (surface of revolution). Az egyes körök mentén elhelyezett
csúcspontokból a program négyszögeket, majd azokat két-két háromszögre bontva lapokat
(face-eket) generál.

A \code{PhongMaterial} anyagmodell a fény és a felszín kölcsönhatását a diffúz
visszaverődés alapján közelíti. Az alkalmazás egy környezeti fényforrást
(\textit{ambient light}) és egy pontfényforrást (\textit{point light}) tartalmaz,
amelyek egymástól függetlenül ki- és bekapcsolhatók.

Az objektumorientált tervezés szempontjából a program az MVC (Model-View-Controller)
mintát követi: a modell tárolja az adatokat, a nézet megjeleníti azokat, a vezérlő
pedig összehangolja a kettőt. A JavaFX property-rendszere lehetővé teszi, hogy az
adatmodell változásai automatikusan tükröződjenek a felhasználói felületen.

% ============================================================
\section{Tervezési dokumentáció}
% ============================================================

\subsection{Fejlesztői környezet}

A program fejlesztéséhez az alábbi eszközök kerültek felhasználásra:

\begin{table}[H]
  \centering
  \begin{tabularx}{\textwidth}{llX}
    \toprule
    \textbf{Eszköz}   & \textbf{Verzió} & \textbf{Felhasználás}                          \\
    \midrule
    Java (JDK)        & 17              & Fejlesztési és futtatási platform               \\
    JavaFX            & 21.0.4          & 3D grafika, kezelőfelület                       \\
    Maven             & 3.9             & Projektmenedzsment, fordítás, függőségkezelés   \\
    IntelliJ IDEA     & 2024.x          & Integrált fejlesztőkörnyezet                    \\
    \bottomrule
  \end{tabularx}
  \caption{A fejlesztéshez használt eszközök}
\end{table}

A \code{pom.xml} konfigurációs fájl tartalmazza a szükséges függőségeket és a
\code{javafx-maven-plugin} bővítményt, amellyel az alkalmazás közvetlenül futtatható
Maven parancsból.

A forráskód a következő főbb JavaFX-technikákat alkalmazza:

\begin{itemize}
  \item \textbf{Property binding:} a csúszkák értéke kötve van a modell tulajdonságaihoz,
        így minden változás automatikusan kiváltja a háló újragenerálását.
  \item \code{TriangleMesh} \textbf{és} \code{MeshView}: egyedi háromszöghálók
        létrehozása és megjelenítése.
  \item \code{PhongMaterial}: Phong-árnyalás az egyes alkatrészekre.
  \item \code{PerspectiveCamera}: perspektív vetítés a 3D jelenethez.
  \item \code{SubScene}: a 3D jelenet és a 2D kezelőfelület elkülönített megjelenítése.
\end{itemize}

\subsection{Csomagstruktúra és osztálydiagram}

A projekt Maven-szabványos könyvtárstruktúrát követ, a forrásfájlok a
\code{src/main/java/com/example/viewer/} mappában találhatók.

\begin{lstlisting}[style=plain, caption={A projekt csomagstruktúrája}]
com.example.viewer
+-- MainApp.java
+-- controller/
|   +-- MainController.java
+-- model/
|   +-- VaseParameters.java
+-- view/
|   +-- ControlPanel.java
|   +-- HelpDialog.java
+-- geometry/
|   +-- VaseMeshGenerator.java
|   +-- SpoutMeshGenerator.java
|   +-- HandleMeshGenerator.java
|   +-- LidDomeMeshGenerator.java
|   +-- LidKnobMeshGenerator.java
+-- utils/
    +-- MathUtils.java
    +-- SceneExporter.java
\end{lstlisting}

\textbf{Osztályok és felelősségeik:}

\textbf{\code{MainApp}} az alkalmazás belépési pontja. Kiterjeszti a JavaFX
\code{Application} osztályt, és a \code{start()} metódusban átadja a vezérlést a
\code{MainController} osztálynak.

\textbf{\code{VaseParameters}} tárolja az összes paramétert JavaFX property objektumok
formájában. Ezek közvetlenül köthetők csúszkákhoz és egyéb vezérlőkhöz. A
\code{handlePos} property felülírja a \code{set()} metódust, hogy a fogantyú pozíciója
mindig a váza magasságán belül maradjon (kényszerfeltétel).

\textbf{\code{MainController}} felügyeli a 3D jelenetet, az egérvezérlést és az összes
\code{MeshView} objektumot. Létrehozza és inicializálja az összes alkatrészt, kezeli a
menüpontokat, és összehangolja a fényforrások állapotát.

\textbf{\code{ControlPanel}} statikus gyártó osztály. A \code{createControlPanel()}
metódus az összes csúszkát, színválasztót és jelölőnégyzetet összegyűjti egy
görgetőpanelba. A csúszka-értékek megváltozásakor minden alkatrész hálója
újragenerálódik.

\textbf{\code{HelpDialog}} egy modális párbeszédablak, amely a program kezelési
útmutatóját tartalmazza.

\textbf{A \code{geometry} csomag} osztályai mind egyetlen statikus metódust tartalmaznak,
amely visszaadja az adott alkatrész \code{TriangleMesh} objektumát. Ez a kialakítás
biztosítja, hogy az egyes geometriagenerátorok egymástól függetlenül fejleszthetők és
tesztelhetők legyenek.

\textbf{\code{MathUtils}} matematikai segédmetódusokat tartalmaz: a váza profil-sugarának
kiszámítását, a \code{smoothStep} interpolációs függvényt és a kiöntő
súlymaszk-számítását.

\textbf{\code{SceneExporter}} az exportálási és importálási logikát valósítja meg.
Külső könyvtár nélkül, kézzel írott JSON-sorosítással dolgozik.

\kephely[10cm]{UML osztálydiagram}

\subsection{Algoritmusok leírása}

\subsubsection{Váza profil és forgástest-generálás}

A váza külső sugarát az alábbi képlet adja meg minden $t$ paraméterértékre
(ahol $t \in [0,\,1]$ a relatív magasság):

\begin{equation}
  r(t) \;=\; r_0
           + b \cdot \sin(\pi t)
           + 9 \cdot \sin(2{,}3\pi t + 0{,}4)
           - n \cdot t^{1{,}4}
\end{equation}

ahol $r_0$ az alapsugár (\code{baseRadius}), $b$ a hasasodás mértéke
(\code{bellyAmount}), $n$ a nyakszűkület (\code{neckTaper}).

A $\sin(\pi t)$ tag biztosítja a hasasodást: az alján és tetején kisebb, középen
nagyobb a sugár. A második szinuszos tag finomabb hullámos felületet ad.
A $n \cdot t^{1{,}4}$ kifejezés a felső rész fokozatos szűkülését modellezi.

A forgástest-generálás lépései:
\begin{enumerate}
  \item A program 64 vízszintes gyűrűt számol ki a magasság mentén.
  \item Minden gyűrűn 72 (alapértelmezés szerint) pontot helyez el egyenlő
        szögtávolságra.
  \item A szomszédos gyűrűk pontjait két háromszöggel köti össze
        (négyszögfelosztás).
  \item Az alsó laphoz középpontot vesz fel, amelyből háromszögek mutatnak kifelé.
  \item A belső falat ugyanígy generálja, de a sugárból kivonja a falvastagságot;
        a felső és alsó perem összekötésével zárt testet kap.
\end{enumerate}

\subsubsection{Kiöntő generálása}

A kiöntő nem külön testként, hanem a váza testének deformációjaként jelenik meg.
Minden csúcspont esetén kiszámít egy \code{spoutWeight} értéket ($\in [0,1]$), amely
megmondja, hogy az adott pont mennyire vesz részt a kiöntőben.

\begin{lstlisting}[style=java, caption={A kiöntő súlymaszk-számítása (MathUtils.java)}]
float topMask     = smoothStep(0.68f, 1.0f, t);        // csak a felso 32%
float widthRad    = toRadians(spoutWidth * 0.5);        // szogszélesség
float angularMask = exp(-(centeredAngle / widthRad)^2); // Gauss-ablak
spoutWeight       = topMask * angularMask;
\end{lstlisting}

A \code{smoothStep} függvény: $s(t) = t^{2}(3 - 2t)$, amely derivált-folytonos
átmenetet biztosít. A kiöntő csúcspontja eltolódik kifelé
($r \mathrel{+}= \text{spoutLength} \cdot w$) és felfelé
($y \mathrel{+}= \text{spoutLift} \cdot w$). Ezáltal a kiöntő sima átmenettel olvad
bele a váza testébe.

\subsubsection{Fogantyú generálása}

A fogantyú egy félkör ívből kihúzott tóruszcikk. Az ív 33 pont mentén halad
(32 szegmens), minden pontban egy 9 csúcsból álló körzetes keresztmetszettel.
A háló $32 \times 8$ négyszöglapot tartalmaz.

A fogantyú pozíciójának korlátja automatikusan érvényesül: a
\code{VaseParameters.handlePos} property \code{set()} metódusa mindig a váza
magasságán belülre szorítja az értéket, a \code{ControlPanel.enforceHandleInsideBounds()}
metódus pedig a fogantyú méretét is korlátozza, hogy ne lógjon ki a teáskanna keretéből.

\subsubsection{Exportálás és importálás}

Az exportálás kézzel írott JSON formátumot alkalmaz, külső könyvtár nélkül.
Az importálás reguláris kifejezésekkel elemzi a fájlt, és ha valamely kulcs hiányzik,
az aktuális értéket tartja meg (graceful degradation). A fájlkezelés szabványos Java
I/O osztályokkal (\code{FileWriter}, \code{BufferedReader}) történik.

% ============================================================
\section{Felhasználói dokumentáció}
% ============================================================

\subsection{Telepítés}

Az alkalmazás futtatásához az alábbiak szükségesek:

\begin{itemize}
  \item \textbf{Java Development Kit (JDK) 17 vagy újabb.} Letölthető az
        \href{https://adoptium.net}{\texttt{adoptium.net}} oldalról. Telepítés után a
        \code{java -version} parancsnak legalább \code{17.x.x} verziószámot kell
        mutatnia.
  \item \textbf{Apache Maven 3.6 vagy újabb.} Letölthető a
        \href{https://maven.apache.org}{\texttt{maven.apache.org}} oldalról. A
        \code{mvn -version} paranccsal ellenőrizhető.
\end{itemize}

A JavaFX könyvtárakat a Maven automatikusan letölti az első fordítás során, külön
telepítés nem szükséges. Az alkalmazás forráskódját tartalmazza a mellékelt ZIP
archívum, amelyet tetszőleges mappába kell kicsomagolni.

\subsection{Futtatás}

A program terminálból (Windows: Command Prompt vagy PowerShell) indítható el.
Először a projekt mappájába kell navigálni:

\begin{lstlisting}[style=shell]
cd C:\[elérési út]\javafx-3d-object-viewer
\end{lstlisting}

Ezt követően a fordítás és futtatás egyetlen paranccsal elvégezhető:

\begin{lstlisting}[style=shell]
mvn javafx:run
\end{lstlisting}

Az első futtatásnál a Maven letölti a szükséges függőségeket, ez néhány percet vehet
igénybe aktív internetkapcsolat esetén. Ha csak fordítani szeretnénk a futtatás nélkül:

\begin{lstlisting}[style=shell]
mvn compile
\end{lstlisting}

\subsection{Menüstruktúra és kezelőfelület}

\subsubsection*{Főablak}

Az alkalmazás ablaka két fő részre oszlik:

\begin{itemize}
  \item \textbf{Bal oldali vezérlőpanel:} itt találhatók a csúszkák, a színválasztók
        és a fény-kapcsolók. A panel görgethető, ha a tartalom nem fér el.
  \item \textbf{Jobb oldali 3D nézet:} a teáskanna interaktívan forgatható
        és nagyítható.
\end{itemize}

\kephely[7cm]{A főablak elrendezése -- vezérlőpanel és 3D nézet}

\subsubsection*{Fájl menü}

\begin{table}[H]
  \centering
  \begin{tabularx}{\textwidth}{lX}
    \toprule
    \textbf{Menüpont}   & \textbf{Funkció}                                                   \\
    \midrule
    Új                  & Véletlenszerű paraméterekkel generál egy új kannaformát            \\
    Törlés              & Az összes paramétert visszaállítja az alapértékekre                \\
    Exportálás\ldots    & Fájlmentési párbeszédet nyit; a paramétereket JSON fájlba menti    \\
    Importálás\ldots    & Fájlmegnyitási párbeszédet nyit; JSON fájlból tölti be a paramétereket \\
    \bottomrule
  \end{tabularx}
  \caption{Fájl menü elemei}
\end{table}

\kephely[4cm]{A Fájl menü nyitott állapotban}

\subsubsection*{Súgó menü}

A \textit{Súgó megjelenítése} menüpont egy modális ablakot nyit meg, amely
összefoglalja az egérvezérlés, a csúszkák, a fájlkezelés és az alkatrészek
kezelésének módját.

\kephely[5cm]{A súgó párbeszédablak}

\subsubsection*{Vezérlőpanel részletei}

\textbf{Szín beállítások:} Négy \code{ColorPicker} vezérlő áll rendelkezésre a test,
a fogantyú, a tető félgömb és a tető fogógombja számára. A szín megváltoztatása
azonnal megjelenik a 3D nézetben.

\textbf{Fények:} Két jelölőnégyzet (\code{CheckBox}) kapcsolja be és ki a környezeti
fényt és a pontfényforrást.

\kephely[6cm]{A vezérlőpanel -- szín és fény beállítások}

\bigskip
\textbf{Váza paraméterek:}

\begin{table}[H]
  \centering
  \begin{tabularx}{\textwidth}{llX}
    \toprule
    \textbf{Csúszka} & \textbf{Tartomány} & \textbf{Leírás}                                         \\
    \midrule
    Magasság         & 170 -- 280         & A váza teljes magassága                                  \\
    Falvastagság     & 4 -- 10            & A váza falának vastagsága                                \\
    Hasasodás        & 6 -- 46            & Mennyire pocakos a váza középső része                    \\
    Nyakszűkület     & 0 -- 22            & Mennyire szűkül a váza fölfelé haladva                   \\
    Alsó sugár       & 30 -- 62           & A váza alapsugarának mérete                              \\
    Felbontás        & 24 -- 128          & A radiális szegmensek száma; nagyobb érték simább felületet ad \\
    \bottomrule
  \end{tabularx}
  \caption{Váza alapparaméterek}
\end{table}

\textbf{Kiöntő paraméterek:}

\begin{table}[H]
  \centering
  \begin{tabularx}{\textwidth}{llX}
    \toprule
    \textbf{Csúszka}   & \textbf{Tartomány} & \textbf{Leírás}                          \\
    \midrule
    Kiöntő hossza      & 4 -- 30            & Mennyire áll ki a kiöntő                 \\
    Kiöntő szélessége  & 15 -- 120          & A kiöntő szögszélessége fokban           \\
    Ajak emelés        & $-20$ -- 18        & A kiöntő ajkának függőleges eltolása     \\
    \bottomrule
  \end{tabularx}
  \caption{Kiöntő paraméterek}
\end{table}

\textbf{Tető paraméterek:}

\begin{table}[H]
  \centering
  \begin{tabularx}{\textwidth}{llX}
    \toprule
    \textbf{Csúszka} & \textbf{Tartomány} & \textbf{Leírás}                            \\
    \midrule
    Tető magasság    & 20 -- 56           & A tetőfélgömb magassága                     \\
    Fogó magasság    & 8 -- 42            & A tető fogógombjának magassága              \\
    Fogó sugár       & 10 -- 20           & A fogógomb alapsugarának mérete             \\
    \bottomrule
  \end{tabularx}
  \caption{Tető paraméterek}
\end{table}

\textbf{Fogantyú paraméterek:}

\begin{table}[H]
  \centering
  \begin{tabularx}{\textwidth}{llX}
    \toprule
    \textbf{Csúszka} & \textbf{Tartomány} & \textbf{Leírás}                                   \\
    \midrule
    Fogó méret       & 20 -- 80           & A fogantyú ívsugarának mérete                      \\
    Fogó pozíció     & $-140$ -- 140      & A fogantyú függőleges helyzete a vázán             \\
    Fogó vastagság   & 4 -- 16            & A fogantyú keresztmetszetének átmérője             \\
    \bottomrule
  \end{tabularx}
  \caption{Fogantyú paraméterek}
\end{table}

\subsubsection*{3D nézet vezérlése}

\begin{table}[H]
  \centering
  \begin{tabularx}{\textwidth}{lX}
    \toprule
    \textbf{Mozdulat}        & \textbf{Hatás}                                       \\
    \midrule
    Bal egérgomb + húzás     & A modell forgatása X és Y tengely körül              \\
    Egérgörgő felfelé        & Közelítés (zoom in)                                  \\
    Egérgörgő lefelé         & Távolítás (zoom out)                                 \\
    \bottomrule
  \end{tabularx}
  \caption{3D nézet egérvezérlése}
\end{table}

A kamera távolságának határa 200 és 2000 egység között van rögzítve, hogy a modell
ne tűnhessen el a látótérből.

\kephely[6.5cm]{A 3D nézet -- forgatott és nagyított teáskanna}

\subsection{Hibakezelés}

\textbf{Importálási hiba:} Ha a kiválasztott JSON fájl olvasása közben hiba lép fel
(pl. sérült fájl, hiányzó kulcsok), az alkalmazás hibaüzenetes párbeszédablakot jelenít
meg, és megtartja az aktuálisan érvényes paramétereket. Részlegesen érvényes fájlok
esetén az olvasható értékek betöltődnek, a hiányzók helyett az előző értékek maradnak.

\textbf{Exportálási hiba:} Ha az exportálás során írási hiba lép fel (pl. nincs írási
jog a kiválasztott mappában), a program szintén hibaüzenetet jelenít meg.

\textbf{Fogantyú korlátok:} A fogantyú pozíciójára és méretére automatikus
kényszerfeltétel vonatkozik. Ha a magasság csúszkájának megváltoztatása után a fogantyú
kilógna a vázáról, a program automatikusan a megengedett tartományba korrigálja az
értéket. Erről értesítés nem jelenik meg, mivel a korrekció folyamatos és a felhasználó
számára láthatatlan marad.

\textbf{Ismeretlen JavaFX hiba:} Ha a Java vagy a JavaFX nincs megfelelően telepítve,
a program terminálban hibaüzenetet ír ki. Ilyen esetben ellenőrizni kell a JDK és a
Maven verzióját.

\kephely[4cm]{Hibaüzenet párbeszédablak -- importálási hiba példa}

% ============================================================
\section{Irodalomjegyzék}
% ============================================================

\begin{enumerate}
  \item Sharan, K.: \textit{Learn JavaFX 17 -- Building User Experience and Interfaces
        with Java}. Apress, 2022.

  \item Bloch, J.: \textit{Effective Java}. 3.~kiadás. Addison-Wesley, 2018.

  \item Oracle Corporation: \textit{JavaFX API Documentation}. [online] Elérhető:
        \url{https://openjfx.io/javadoc/21/} [Letöltve: 2025. január]

  \item OpenJFX Contributors: \textit{OpenJFX Getting Started Guide}. [online]
        Elérhető: \url{https://openjfx.io/openjfx-docs/}
        [Letöltve: 2025. január]

  \item Apache Software Foundation: \textit{Maven Getting Started Guide}. [online]
        Elérhető: \url{https://maven.apache.org/guides/getting-started/}
        [Letöltve: 2025. február]

  \item Hughes, J.~F., Van Dam, A., McGuire, M. et al.:
        \textit{Computer Graphics: Principles and Practice}. 3.~kiadás.
        Addison-Wesley, 2014.
\end{enumerate}

% ============================================================
\section{Melléklet}
% ============================================================

\subsection{Forráskód}

A forráskód a projekt \code{src/main/java/com/example/viewer/} könyvtárában
található. A következő osztályok teljes kódja kerül közlésre; a geometriageneráló
osztályok (\code{geometry} csomag) a mellékelt ZIP archívumban olvashatók.

\subsubsection*{MainApp.java}

\begin{lstlisting}[style=java, caption={MainApp.java -- belépési pont}]
package com.example.viewer;

import com.example.viewer.controller.MainController;
import javafx.application.Application;
import javafx.stage.Stage;

public class MainApp extends Application {

    private MainController controller;

    @Override
    public void start(Stage stage) {
        controller = new MainController();
        controller.start(stage);
    }

    public static void main(String[] args) {
        launch(args);
    }
}
\end{lstlisting}

\subsubsection*{model/VaseParameters.java}

\begin{lstlisting}[style=java, caption={VaseParameters.java -- adatmodell}]
package com.example.viewer.model;

import javafx.beans.property.DoubleProperty;
import javafx.beans.property.IntegerProperty;
import javafx.beans.property.SimpleDoubleProperty;
import javafx.beans.property.SimpleIntegerProperty;
import javafx.scene.paint.Color;

public class VaseParameters {
    public final DoubleProperty height          = new SimpleDoubleProperty(260);
    public final DoubleProperty wallThickness   = new SimpleDoubleProperty(10);
    public final DoubleProperty bellyAmount     = new SimpleDoubleProperty(22);
    public final DoubleProperty neckTaper       = new SimpleDoubleProperty(4);
    public final DoubleProperty baseRadius      = new SimpleDoubleProperty(38);
    public final IntegerProperty radialSegments = new SimpleIntegerProperty(72);

    public final DoubleProperty spoutLength     = new SimpleDoubleProperty(12);
    public final DoubleProperty spoutWidth      = new SimpleDoubleProperty(70);
    public final DoubleProperty spoutLift       = new SimpleDoubleProperty(7);

    public final DoubleProperty lidHeight       = new SimpleDoubleProperty(30);
    public final DoubleProperty knobHeight      = new SimpleDoubleProperty(20);
    public final DoubleProperty knobRadius      = new SimpleDoubleProperty(12);

    public final DoubleProperty handleSize      = new SimpleDoubleProperty(40);
    public final DoubleProperty handleThickness = new SimpleDoubleProperty(8);

    public final DoubleProperty handlePos = new SimpleDoubleProperty(0) {
        @Override
        public void set(double value) {
            double min = -height.get() / 2;
            double max =  height.get() / 2;
            super.set(Math.max(min, Math.min(max, value)));
        }
    };

    public final javafx.beans.property.ObjectProperty<Color> bodyColor =
        new javafx.beans.property.SimpleObjectProperty<>(Color.rgb(210, 130, 80));
    public final javafx.beans.property.ObjectProperty<Color> handleColor =
        new javafx.beans.property.SimpleObjectProperty<>(Color.rgb(180, 100, 60));
    public final javafx.beans.property.ObjectProperty<Color> lidDomeColor =
        new javafx.beans.property.SimpleObjectProperty<>(Color.rgb(196, 116, 72));
    public final javafx.beans.property.ObjectProperty<Color> lidKnobColor =
        new javafx.beans.property.SimpleObjectProperty<>(Color.rgb(220, 140, 80));

    public final javafx.beans.property.BooleanProperty ambientLightEnabled =
        new javafx.beans.property.SimpleBooleanProperty(true);
    public final javafx.beans.property.BooleanProperty pointLightEnabled =
        new javafx.beans.property.SimpleBooleanProperty(true);
}
\end{lstlisting}

\subsubsection*{utils/MathUtils.java}

\begin{lstlisting}[style=java, caption={MathUtils.java -- matematikai segédfüggvények}]
package com.example.viewer.utils;

public class MathUtils {

    public static float smoothStep(float edge0, float edge1, float x) {
        float t = (x - edge0) / (edge1 - edge0);
        t = Math.max(0f, Math.min(1f, t));
        return t * t * (3f - 2f * t);
    }

    public static float findMaxVaseRadius(
            com.example.viewer.model.VaseParameters params) {
        float maxRadius = 0f;
        for (int i = 0; i <= 120; i++) {
            float t = (float) i / 120f;
            maxRadius = Math.max(maxRadius, vaseProfileRadius(t, params));
        }
        return maxRadius;
    }

    public static float vaseProfileRadius(
            float t, com.example.viewer.model.VaseParameters params) {
        return (float) params.baseRadius.get()
             + (float) params.bellyAmount.get() * (float) Math.sin(Math.PI * t)
             + 9f * (float) Math.sin(2.3 * Math.PI * t + 0.4)
             - (float) params.neckTaper.get() * (float) Math.pow(t, 1.4);
    }

    public static float computeSpoutWeight(
            float t, double angle,
            com.example.viewer.model.VaseParameters params) {
        float topMask    = smoothStep(0.68f, 1.0f, t);
        float widthRad   = (float) Math.toRadians(params.spoutWidth.get() * 0.5);
        widthRad         = Math.max(0.12f, widthRad);
        float centered   = (float) Math.atan2(
            Math.sin(angle - Math.PI), Math.cos(angle - Math.PI));
        float angularMask = (float) Math.exp(
            -Math.pow(centered / widthRad, 2.0));
        return topMask * angularMask;
    }
}
\end{lstlisting}

\subsection{Értékelő lap}

\begin{table}[H]
  \centering
  \begin{tabular}{lcc}
    \toprule
    \textbf{Szempont}       & \textbf{Max. pont} & \textbf{Kapott pont} \\
    \midrule
    Dokumentáció            & 5                  &                      \\
    Program működés         & 5                  &                      \\
    Megvalósítás minősége   & 5                  &                      \\
    Bemutató                & 5                  &                      \\
    \midrule
    \textbf{Összesen}       & \textbf{20}        &                      \\
    \bottomrule
  \end{tabular}
  \caption{Értékelő lap}
\end{table}

\begin{table}[H]
  \centering
  \begin{tabular}{cc}
    \toprule
    \textbf{Pont} & \textbf{Jegy}    \\
    \midrule
    0 -- 7        & 1 (elégtelen)    \\
    8 -- 10       & 2 (elégséges)    \\
    11 -- 13      & 3 (közepes)      \\
    14 -- 16      & 4 (jó)           \\
    17 -- 20      & 5 (jeles)        \\
    \bottomrule
  \end{tabular}
  \caption{Ponthatárok és jegyek}
\end{table}

\bigskip
\noindent\textbf{Elért jegy:}\quad \underline{\hspace{3.5cm}}

\bigskip
\noindent\textbf{Megjegyzés:}\quad \underline{\hspace{9cm}}

\bigskip
\noindent\textbf{Aláírás:}\quad \underline{\hspace{5cm}}

\end{document}
